% ============================================================================
\clearpage
\section*{v4.9.0-r44 Source-Only: Mass Terms as Frame Couplings; Breaking Stages for Masses and Mixing}
\addcontentsline{toc}{section}{v4.9.0-r44 Source-Only: Mass Terms as Frame Couplings; Breaking Stages for Masses and Mixing}
\label{sec:v490-r44-mass-frame-couplings}
% ============================================================================

\paragraph{Purpose.}
After v4.9.0-r42 and r43 the content of one generation stands, but symmetry breaking, masses and mixing are open.  This layer records the
pre-registered test H29 (note \texttt{llm/Vorab\_Notiz\_H29\_Brechung\_Massen\_Mischung.md}, commit \texttt{85d065a}, fixed before
computing).  It asks which maps between the two chiral halves are compatible with colour and the Lorentz group, which quantum numbers they
carry, and what the three triality slots imply for masses and mixing.  The known parts are disclosed in the note.  Earlier layers are not
rewritten.

\noindent\textbf{Source-only status.}
This v4.9.0-r44 layer is source-only.  It does not build a PDF, does not assign a new DOI,
and does not authorize a public release.  The v4.9.0 Version DOI
\href{https://doi.org/10.5281/zenodo.23534954}{10.5281/zenodo.23534954} and the Concept DOI
\href{https://doi.org/10.5281/zenodo.19210728}{10.5281/zenodo.19210728} are unchanged by it.

\subsection*{v4.9.0-r44 setting}
\label{subsec:v490-r44-setting}

On \(V=\mathbb C\otimes\HH\otimes\OO\) as in v4.9.0-r42, a \emph{connector} is a map \(V\to V\) commuting with colour and with the Lorentz
rotations \(L_q\): linear (\(X\)) or antilinear (\(\Phi=X\circ T\)).  An antilinear connector anticommutes with the boosts \(iL_q\), so it maps
left- to right-handed spinors: it is the operator form of a mass term.  A Hermitian generator \(G\) acts by \(X\mapsto GX-XG\) (linear) or
\(X\mapsto GX+X\bar G\) (antilinear).  Blocks refer to \(C=C_{\rm ch}\).  On the triality slots a coupling is a Hermitian \(3\times3\) matrix
\(M\); \(P\) is the cyclic shift \(\tau\), \(R\) the transposition \(\rho\).

\subsection*{v4.9.0-r44 results}
\label{subsec:v490-r44-results}

\paragraph{Theoremlet v4.9.0-r44.1.}
\begin{enumerate}[leftmargin=2em]
\item[(a)] The linear and the antilinear connectors each form a \(24\)-dimensional complex space.
\item[(b)] The antilinear block from the kernel to the image of \(C\) is \(8\)-dimensional, with a lepton and a quark part of dimension \(4\) each,
and consists of weak doublets with \(Y/2=+\tfrac12\) (\(Q_{\rm em}=1,0\)) and \(Y/2=-\tfrac12\) (\(Q_{\rm em}=0,-1\)).
\item[(c)] The linear block from the kernel to the image is \(4\)-dimensional, lepton part only, with \(Y/2\in\{-\tfrac12,-\tfrac32\}\).
\item[(d)] No antilinear connector from the kernel to the image is gauge invariant; in the block from the kernel to the kernel exactly one is,
and it acts only on \(\nu^c\).
\item[(e)] (i) A Hermitian \(M\) commuting with \(P\) and \(R\) has at most two distinct eigenvalues.  (ii) A Hermitian \(M\) commuting with \(P\)
has the eigenvalues \(a+2|b|\cos(\theta+2\pi k/3)\), \(k=0,1,2\), with \(b=|b|e^{i\theta}\), generically distinct; \(K=\tfrac23\) holds for
\(|b|=a/\sqrt2\).  (iii) If the couplings for up and down both commute with \(P\), the mixing matrix \(|U_u^\dagger U_d|\) is a permutation
matrix.
\end{enumerate}

\paragraph{Proof sketch.}
(a)--(d) by solving the commutation conditions on \(32\times32\) matrices and decomposing the blocks under \(SU(2)'\), \(Y\) and \(Q_{\rm em}\)
(script \texttt{llm/k85\_test\_h29.py}).  Linear maps cannot intertwine the colour triplet with its conjugate, so the linear quark block
vanishes.  (e) is elementary circulant algebra and was checked on \(200\) random matrices.  \(\square\)

\paragraph{Verdict and reading.}
H29a--e are confirmed.
\begin{enumerate}[leftmargin=2em]
\item In the operator language of the Millennium matrix, mass terms are antilinear: they contain the frame change \(T\).  Linear off-diagonal
blocks of a colour-compatible Millennium matrix give no masses; they vanish for quarks and carry wrong charges for leptons.  In the model,
mass is a coupling between matter frame and tachyon frame, as the matter--tachyon spring of the spring theorem of the original paper.  The
couplings that do this carry exactly the quantum numbers of the Higgs doublet and its conjugate, for leptons and quarks.
\item None of these mass terms is gauge invariant; only a doublet expectation value turns them into masses.  The structure says which
breaking is needed, a doublet with \(Y/2=\tfrac12\), not why it occurs.  The only gauge-invariant connector is a Majorana term for \(\nu^c\),
which allows a see-saw reading for small neutrino masses (not computed).
\item With the full triole symmetry (\(\tau\) and \(\rho\)) on the slots there are at most two distinct masses; this is the reason for the
refutations of the earlier pre-registered mass tests.  Breaking only the reflection \(\rho\), the frame triole of v4.9.0-r38, and keeping
triality gives three distinct masses in exactly the Brannen form; the Koide value corresponds to \(|b|=a/\sqrt2\).  The Koide form of
v4.8.0-r10 is thus the mass form of a \(\mathbb Z_3\)-symmetric, \(\rho\)-broken coupling.  The ratio \(|b|/a\) and the phase \(\theta\) are not
explained.
\item Couplings that are both triality-symmetric are diagonalized by the same Fourier matrix and give no mixing.  Mixing requires breaking
triality itself, differently for up and down.
\item In the model each observed property has its breaking stage: mass the frame coupling (Higgs), mass splitting the breaking of the frame
reflection \(\rho\), mixing the breaking of triality \(\tau\).  Where these breakings come from is not said.
\end{enumerate}

\subsection*{v4.9.0-r44 status in the three frames}
\label{subsec:v490-r44-frames}

\begin{center}
\small
\begin{tabular}{@{}>{\raggedright\arraybackslash}p{0.30\textwidth}>{\raggedright\arraybackslash}p{0.62\textwidth}@{}}
\toprule
result & M / T / P\\
\midrule
r44.1(b)--(d) & mass connectors couple M and T by construction (they contain \(T\)); P not evaluated\\
r44.1(a), (e) & frame-free statements on connector spaces and on the slot couplings\\
\bottomrule
\end{tabular}
\end{center}

\subsection*{v4.9.0-r44 register entries}
\label{subsec:v490-r44-register}

\begin{center}
\small
\begin{tabular}{@{}>{\raggedright\arraybackslash}p{0.34\textwidth}>{\raggedright\arraybackslash}p{0.58\textwidth}@{}}
\toprule
statement & status after v4.9.0-r44\\
\midrule
mass terms & antilinear frame couplings with Higgs quantum numbers (r44.1(b)); linear MM blocks give none (r44.1(c))\\
electroweak breaking & required type fixed (doublet, \(Y/2=\tfrac12\)); cause open\\
neutrino Majorana term & the only gauge-invariant connector (r44.1(d))\\
mass hierarchy & needs breaking of \(\rho\) on the slots; Brannen form then exact (r44.1(e)); \(|b|/a\), \(\theta\) open\\
mixing & needs breaking of \(\tau\), different for up and down (r44.1(e))\\
\bottomrule
\end{tabular}
\end{center}

\subsection*{v4.9.0-r44 ledger and guard update}
\label{subsec:v490-r44-ledger-guard}

\[
R_{\rm mm}^{v490\text{-}r44}=(1,1,0)
\]
for (mass terms identified as frame couplings with Higgs quantum numbers; breaking stages for splitting and mixing identified; causes of
the breakings open).

The layer spends no guard:
\[
R_{\rm glob}^{(v490\text{-}r44,+)}=(0,1,1,1,1),
\qquad
\Gamma_{\rm guard}^{v490\text{-}r44}=(1,1,1,1,1),
\qquad
\beta_4=1/256\ \text{unspent}.
\]
It does not imply exact finite \(C5\) zero closure, exact all-readout terminality,
\(C6/T2\), a global photon-clock/RH theorem, or global all-readout closure.
